%% kapa: exam=1 round=33 year=2022 period=1 form=A subject=real_estate subject_ko=부동산학원론 seq=29 number=109
%% source: Q-Net 감정평가사 기출문제 (cst00309) artlSeq=5209997 file=raw/1008/5209997/2022년도 제33회 감정평가사 1차시험 1교시(민법,경제학,부동산학).pdf
%% conversion: kordoc 4.24.51 -> tools/gen1.py (auto); review=proofread-sonnet; flags=none
\begin{question}{109}
부동산 중개계약에 관한 설명으로 옳지 않은 것은?
\begin{choices}
  \item 순가중개계약에서는 매도자가 개업공인중개사에게 제시한 가격을 초과해 거래가 이루어진 경우 그 초과액을 매도자와 개업공인중개사가 나누어 갖는다.
  \item 일반중개계약에서는 의뢰인이 다수의 개업공인중개사에게 동등한 기회로 거래를 의뢰한다.
  \item 공인중개사법령상 당사자간에 다른 약정이 없는 경우 전속중개계약의 유효기간은 3월로 한다.
  \item 공동중개계약에서는 부동산거래정보망 등을 통하여 다수의 개업공인중개사가 상호 협동하여 공동으로 거래를 촉진한다.
  \item 독점중개계약에서는 의뢰인이 직접 거래를 성사시킨 경우에도 중개보수 청구권이 발생한다.
\end{choices}
\end{question}
\answer{1}
